\subsection{Математическое ожидание}

Вторая по важности характеристика случайной величины -- математическое ожидание. Именно для неё мы ввели интеграл Лебега:

\begin{definition}[Математическое ожидание]
Математическим ожиданием случайной величины $\xi$ на вероятностном пространстве $(\Omega, \mathfrak{S}, \mathbb{P})$ называется следующая величина:
\[
\mathbb{E}\xi = \int_{\Omega} \xi(\omega) \, \mathbb{P}(\mathrm{d}\omega).
\]
\end{definition}

\begin{theorem}[Вычисление математического ожидания через распределение]
Математическое ожидание случайной величины $\xi$ с распределением $\mathbb{P}_{\xi}$ равно
\[
\mathbb{E}\xi = \int_{-\infty}^{+\infty} x \, \mathbb{P}_{\xi}(\mathrm{d}x).
\]
\end{theorem}

\textit{Доказательство.} Докажем это утверждение последовательно расширяя класс функций.

\textbf{Шаг 1. Индикаторные функции.} Пусть $\xi(\omega) = \mathbf{1}_A(\omega)$ для некоторого $A \in \mathfrak{S}$. Тогда по определению интеграла Лебега:
\[
\int_{\Omega} \mathbf{1}_A(\omega) \, \mathbb{P}(\mathrm{d}\omega) = \mathbb{P}(A).
\]
С другой стороны, случайная величина $\xi$ принимает только два значения: $1$ (на множестве $A$) и $0$ (на $\Omega \setminus A$). Её распределение $\mathbb{P}_{\xi}$ сосредоточено в точках $0$ и $1$:
\[
\mathbb{P}_{\xi}(\{1\}) = \mathbb{P}(A), \quad \mathbb{P}_{\xi}(\{0\}) = \mathbb{P}(\Omega \setminus A).
\]
Тогда:
\[
\int_{-\infty}^{+\infty} x \, \mathbb{P}_{\xi}(\mathrm{d}x) = 1 \cdot \mathbb{P}_{\xi}(\{1\}) + 0 \cdot \mathbb{P}_{\xi}(\{0\}) = \mathbb{P}(A).
\]
Равенство доказано для индикаторов.

\textbf{Шаг 2. Простые функции.} Пусть $\xi$ -- простая функция, т.е. $\xi(\omega) = \sum_{k=1}^{n} c_k \mathbf{1}_{A_k}(\omega)$, где $A_k \in \mathfrak{S}$ -- попарно непересекающиеся множества, $\bigcup_{k=1}^{n} A_k = \Omega$. По линейности интеграла Лебега и результату Шага 1:
\[
\int_{\Omega} \xi(\omega) \, \mathbb{P}(\mathrm{d}\omega) = \sum_{k=1}^{n} c_k \int_{\Omega} \mathbf{1}_{A_k}(\omega) \, \mathbb{P}(\mathrm{d}\omega) = \sum_{k=1}^{n} c_k \, \mathbb{P}(A_k).
\]
С другой стороны, распределение $\mathbb{P}_{\xi}$ сосредоточено в точках $c_1, \ldots, c_n$, причём $\mathbb{P}_{\xi}(\{c_k\}) = \mathbb{P}(A_k)$. Поэтому:
\[
\int_{-\infty}^{+\infty} x \, \mathbb{P}_{\xi}(\mathrm{d}x) = \sum_{k=1}^{n} c_k \, \mathbb{P}_{\xi}(\{c_k\}) = \sum_{k=1}^{n} c_k \, \mathbb{P}(A_k).
\]
Равенство доказано для простых функций.

\textbf{Шаг 3. Неотрицательные случайные величины.} Пусть $\xi(\omega) \geqslant 0$. По свойствам измеримых функций, существует неубывающая последовательность простых функций $\xi_n(\omega) \nearrow \xi(\omega)$ при $n \to \infty$. По теореме Леви о монотонной сходимости:
\[
\int_{\Omega} \xi(\omega) \, \mathbb{P}(\mathrm{d}\omega) = \lim_{n \to \infty} \int_{\Omega} \xi_n(\omega) \, \mathbb{P}(\mathrm{d}\omega).
\]
По Шагу 2 для каждого $n$ выполняется:
\[
\int_{\Omega} \xi_n(\omega) \, \mathbb{P}(\mathrm{d}\omega) = \int_{-\infty}^{+\infty} x \, \mathbb{P}_{\xi_n}(\mathrm{d}x).
\]
Поскольку $\xi_n \nearrow \xi$, распределения $\mathbb{P}_{\xi_n}$ сходятся к $\mathbb{P}_{\xi}$, и по теореме о монотонной сходимости (применённой уже к интегралу по $\mathbb{R}$):
\[
\lim_{n \to \infty} \int_{-\infty}^{+\infty} x \, \mathbb{P}_{\xi_n}(\mathrm{d}x) = \int_{-\infty}^{+\infty} x \, \mathbb{P}_{\xi}(\mathrm{d}x).
\]
Равенство доказано для неотрицательных случайных величин.

\textbf{Шаг 4. Общий случай.} Произвольную случайную величину $\xi$ представим в виде разности положительной и отрицательной частей: $\xi = \xi^+ - \xi^-$, где $\xi^+ = \max(\xi, 0)$ и $\xi^- = \max(-\xi, 0)$. Обе части неотрицательны, поэтому по Шагу 3 и линейности интеграла:
\[
\int_{\Omega} \xi(\omega) \, \mathbb{P}(\mathrm{d}\omega) = \int_{\Omega} \xi^+(\omega) \, \mathbb{P}(\mathrm{d}\omega) - \int_{\Omega} \xi^-(\omega) \, \mathbb{P}(\mathrm{d}\omega) = \int_{-\infty}^{+\infty} x \, \mathbb{P}_{\xi}(\mathrm{d}x). \hfill \blacksquare
\]

Данную теорему можно было бы оставить и без доказательства, хотя она довольно очевидна, поскольку здесь мы просто делаем замену $x = \xi(\omega)$, а $\mathbb{P}(\Omega_0) = \mathbb{P}_{\xi}(\xi(\Omega_0))$, где $\Omega_0 \in \mathfrak{S}$.

\bigskip

Часто математическое ожидание не имеет четкой интерпретации. Например, как проинтерпретировать математическое ожидание выпавшего на игральном кубике после броска числа (равного $3,5$, рекомендую проверить это самостоятельно)? Я бы сказал, что матожидание - это ответ на вопрос: "Какое число требуется взять, чтобы оно было наиболее близким к любому из возможных значений случайной величины?"

\begin{notice}
Далее мы докажем, что среднее арифметическое случайных величин с одинаковым распределением сходится к своему математическому ожиданию.\\

Важно отметить, что математического ожидания у случайных величин может и не существовать (если интеграл Лебега расходится). В будущем будут некоторые примеры таких величин.
\end{notice}

\subsubsection{Теорема Лебега о мажорируемой сходимости}
\begin{theorem}[Теорема Лебега о мажорируемой сходимости]
Пусть $\{\xi_n\}_{n=1}^{\infty}$ -- последовательность случайных величин. Тогда, если $\xi_n \to \xi$ почти всюду, и пусть существует случайная величина $\eta$ с конечным математическим ожиданием $\forall n \in \mathbb{N}$ $|\xi_n| \leqslant \eta$ почти наверное (т.е. вероятность обратного -- ноль). Тогда существуют матожидания $\mathbb{E}\xi_n$ и $\mathbb{E}\xi$ и при этом
\[
\lim_{n \to \infty} \mathbb{E}\xi_n = \mathbb{E}\xi.
\]
\end{theorem}

Данная теорема -- следствие теоремы Лебега о мажорируемой сходимости.

\begin{notice}
\textbf{Развернутое пояснение к теореме о мажорируемой сходимости.}

Эта теорема отвечает на фундаментальный вопрос: \textit{когда можно поменять местами предел и математическое ожидание (интеграл)?} То есть когда верно равенство
\[
\mathbb{E}\!\left[\lim_{n \to \infty} \xi_n\right] = \lim_{n \to \infty} \mathbb{E}[\xi_n]?
\]

На первый взгляд кажется, что это должно выполняться всегда, если $\xi_n \to \xi$. Однако это \textbf{не так}. Проблема в том, что даже если значения $\xi_n(\omega)$ сходятся к $\xi(\omega)$ для почти каждого $\omega$, их интегралы (матожидания) могут вести себя ``плохо'' -- масса распределения может ``убегать на бесконечность''.

\textbf{Зачем нужна мажоранта $\eta$?}

Условие $|\xi_n| \leqslant \eta$ с $\mathbb{E}\eta < \infty$ гарантирует, что ``хвосты'' распределений $\xi_n$ не разбегаются. Мажоранта $\eta$ как бы ``запирает'' все $\xi_n$ в рамках одной интегрируемой функции, не давая массе уходить в бесконечность.

\textbf{Контрпример без мажоранты.} Рассмотрим $\Omega = (0, 1]$ с равномерной мерой и $\xi_n(\omega) = n \cdot \mathbf{1}_{(0, 1/n]}(\omega)$. Тогда:
\begin{itemize}
    \item Для любого фиксированного $\omega > 0$ при достаточно больших $n$ имеем $\omega > 1/n$, значит $\xi_n(\omega) = 0$. Следовательно, $\xi_n \to 0$ поточечно (и почти всюду).
    \item Но $\mathbb{E}\xi_n = n \cdot \frac{1}{n} = 1$ для всех $n$.
    \item Значит, $\lim_{n \to \infty} \mathbb{E}\xi_n = 1 \neq 0 = \mathbb{E}[0]$.
\end{itemize}
Здесь мажоранты с конечным матожиданием не существует: $\sup_n \xi_n(\omega) = \lfloor 1/\omega \rfloor$, и её интеграл расходится.

\textbf{Практический смысл.} В задачах теории вероятностей эта теорема позволяет обосновывать переход к пределу под знаком матожидания -- например, при доказательстве закона больших чисел или при работе с рядами случайных величин.
\end{notice}

\begin{definition}[Борелевская функция]
\textbf{Борелевской функцией} называется отображение одного топологического пространства в другое, т.е. это такое отображение, в котором прообраз любого открытого множества -- открытое множество.
\end{definition}

\begin{theorem}
Если $f: \mathbb{R}^n \to \mathbb{R}$ -- борелевская функция, а $\xi_1, \ldots, \xi_n$ -- некоторые случайные величины на одном вероятностном пространстве, то $f(\xi_1, \xi_2, \ldots, \xi_n)$ -- случайная величина.
\end{theorem}

Все непрерывные функции -- борелевские. В частности, если $\xi$ и $\eta$ -- случайные величины, то $a\xi + b\eta$, $\xi \cdot \eta$ и $\dfrac{\xi}{\eta}$ (при условии $\eta \neq 0$) -- тоже случайные величины.

\begin{notice}
\textbf{Развернутое пояснение о борелевских функциях.}

\textbf{Зачем вообще нужно это определение?} Вспомним, что случайная величина -- это \textit{измеримая} функция $\xi: \Omega \to \mathbb{R}$. Измеримость означает, что прообраз любого борелевского множества из $\mathbb{R}$ лежит в $\sigma$-алгебре $\mathfrak{S}$ на $\Omega$. Это нужно, чтобы мы могли говорить о вероятности $\mathbb{P}(\xi \in B)$ для любого ``разумного'' множества $B$.

Теперь возникает естественный вопрос: если мы взяли случайные величины $\xi_1, \ldots, \xi_n$ и применили к ним какую-то функцию $f$, останется ли результат $f(\xi_1, \ldots, \xi_n)$ случайной величиной? То есть останется ли он измеримым?

\textbf{Ответ:} да, если $f$ -- борелевская функция. Это следует из того, что композиция измеримых функций измерима.

\textbf{Почему непрерывные функции всегда борелевские?} По определению непрерывности, прообраз любого открытого множества при непрерывном отображении открыт. А любое открытое множество в $\mathbb{R}$ -- борелевское. Следовательно, прообраз борелевского множества при непрерывной функции -- борелевское множество.

\textbf{Практический вывод.} В подавляющем большинстве задач мы работаем с ``хорошими'' функциями: полиномами, экспонентами, логарифмами, тригонометрическими функциями, их композициями. Все они непрерывны (или кусочно-непрерывны), а значит борелевские. Поэтому можно спокойно писать $\xi^2$, $e^{\xi}$, $\sin(\xi)$, $\xi + \eta$ и быть уверенным, что это корректно определённые случайные величины.

\textbf{Нюанс с делением.} Функция $f(x, y) = x/y$ не определена при $y = 0$, поэтому $\xi/\eta$ является случайной величиной только на множестве $\{\eta \neq 0\}$. Если $\mathbb{P}(\eta = 0) = 0$, то это не создаёт проблем -- мы можем доопределить значение на этом множестве произвольно (например, нулём), и измеримость сохранится.
\end{notice}
\subsubsection{Свойства математического ожидания}
\begin{theorem}[Свойства математического ожидания]
Пусть $\xi_1, \xi_2, \ldots, \xi_n$ -- случайные величины, определённые на одном вероятностном пространстве, а $f: \mathbb{R}^n \to \mathbb{R}$ -- борелевская функция.
\begin{enumerate}
    \item $\forall c \in \mathbb{R}$ $\mathbb{E}[\xi_1 + c] = \mathbb{E}\xi_1 + c$;
    \item $\forall a, b \in \mathbb{R}$ $\mathbb{E}[a\xi_1 + b\xi_2] = a\mathbb{E}\xi_1 + b\mathbb{E}\xi_2$;
    \item $\mathbb{E}[f(\xi_1, \xi_2, \ldots, \xi_n)] = \displaystyle\int_{\Omega} f(\xi_1(\omega), \xi_2(\omega), \ldots, \xi_n(\omega)) \, \mathbb{P}(\mathrm{d}\omega)$.
\end{enumerate}
\end{theorem}

\textit{Доказательство.} Третье свойство очевидно из определения математического ожидания и из теоремы 2.2.3. Остальные два свойства -- следствия третьего. Действительно, например, первое свойство:
\[
\mathbb{E}[\xi + c] = \int_{\Omega} (\xi(\omega) + c) \, \mathbb{P}(\mathrm{d}\omega) = \int_{\Omega} \xi(\omega) \, \mathbb{P}(\mathrm{d}\omega) + c \int_{\Omega} \mathbb{P}(\mathrm{d}\omega) = \mathbb{E}\xi_1 + c\,\mathbb{P}(\Omega) = \mathbb{E}\xi_1 + c.
\]
Второе свойство доказывается аналогично с помощью аддитивности интеграла Лебега. $\blacksquare$

\begin{notice}
\textbf{Развернутое пояснение к свойствам математического ожидания.}

\textbf{Свойство 1 (сдвиг).} Матожидание константы равно самой константе: $\mathbb{E}[c] = c$. Это интуитивно понятно -- если величина всегда принимает одно и то же значение $c$, то её ``среднее'' и есть $c$. Формально это следует из того, что $\int_{\Omega} c \, \mathbb{P}(\mathrm{d}\omega) = c \cdot \mathbb{P}(\Omega) = c \cdot 1 = c$.

\textbf{Свойство 2 (линейность).} Это одно из важнейших свойств матожидания. Оно утверждает, что оператор $\mathbb{E}$ линеен:
\[
\mathbb{E}[a\xi + b\eta] = a\,\mathbb{E}\xi + b\,\mathbb{E}\eta.
\]
Обратите внимание: линейность выполняется \textbf{без каких-либо предположений} о независимости $\xi$ и $\eta$! Это отличает матожидание от дисперсии, где $\mathbb{D}[\xi + \eta] = \mathbb{D}\xi + \mathbb{D}\eta$ только при независимости (или некоррелированности).

\textbf{Свойство 3 (замена переменной / теорема о функции от случайной величины).} Это свойство -- прямое обобщение теоремы 2.2.1. Оно говорит, что матожидание функции от случайных величин можно вычислять двумя эквивалентными способами:
\begin{itemize}
    \item \textbf{В пространстве исходов $\Omega$:} интегрировать $f(\xi_1(\omega), \ldots, \xi_n(\omega))$ по мере $\mathbb{P}$ на $\Omega$.
    \item \textbf{В пространстве значений $\mathbb{R}^n$:} интегрировать $f(x_1, \ldots, x_n)$ по совместному распределению $(\xi_1, \ldots, \xi_n)$.
\end{itemize}

\textbf{Практическая важность.} Чаще всего мы работаем вторым способом -- через распределение. Например, если $\xi$ имеет плотность $f_{\xi}$, то
\[
\mathbb{E}[g(\xi)] = \int_{-\infty}^{+\infty} g(x) f_{\xi}(x) \, \mathrm{d}x.
\]
Это позволяет вычислять матожидания функций от случайных величин, \textbf{не находя предварительно распределение} $g(\xi)$. Например, чтобы найти $\mathbb{E}[\xi^2]$, не нужно искать распределение $\xi^2$ -- достаточно проинтегрировать $x^2 f_{\xi}(x)$.
\end{notice}